% SOURCE_AUTHORITY: sga5_fr_workpass.tex, lines 15042--15066 (LF-normalized unit).
% SOURCE_SHA256: 791F4EFFC5E02832D5D77ED03518C8156D6F07E4C8238B03545DB93D883FBB28
\begin{statement}{Lema 3}
Seja $X$ um esquema de tipo finito sobre $\mathbb F_p$, e seja $g:X\to e$ seu morfismo estrutural.
\begin{enumerate}[label=\alph*)]
\item Consideremos um $\Omega$-feixe construtível $F$ sobre $X$, um subesquema fechado $Y$ de $X$ e o aberto complementar $U=X-Y$. Se o teorema for verdadeiro para dois dos pares $(X,F)$, $(Y,F|Y)$, $(U,F|U)$, será verdadeiro para o terceiro.
\item Consideremos uma sequência exata
\[
0\longrightarrow F'\longrightarrow F\longrightarrow F''\longrightarrow0
\]
de $\Omega$-feixes construtíveis sobre $X$; se o teorema for verdadeiro para dois dos pares $(X,F')$, $(X,F)$, $(X,F'')$, será verdadeiro para o terceiro.
\end{enumerate}
\end{statement}

Utilizaremos apenas o enunciado a), mas damos também b) por simetria. As demonstrações de a) e b) são análogas; demonstremos, por exemplo, a). A sequência exata
\[
\cdots\longrightarrow \R^i_!g_U(F|U)\longrightarrow \R^i_!g(F)
\longrightarrow \R^i_!g_Y(F|Y)\longrightarrow \R^{i+1}_!g_U(F|U)\longrightarrow\cdots
\]
e a proposição 1a) (n\textsuperscript{o}~1) dão
\[
\prod_i\bigl(L_{\R^i_!g(F)}\bigr)^{(-1)^i}
=\prod_i\bigl(L_{\R^i_!g_U(F|U)}\,L_{\R^i_!g_Y(F|Y)}\bigr)^{(-1)^i},
\]
e a proposição 1b) (n\textsuperscript{o}~1) permite então concluir.

Os raciocínios dos lemas 1, 2 e 3 podem ser resumidos dizendo que, segundo a proposição 1 do n\textsuperscript{o}~1, a função $L_F$ tem propriedades formais análogas às da cohomologia com suportes próprios.
